% Proposed abstract for the v2 manuscript; approximately 190 words.
% Title options are supplied separately to the integrating author.
\begin{abstract}
Transitions into two symmetry-related absorbing states in two dimensions are closely
associated with generalized voter criticality or with separate symmetry-breaking and
absorption transitions. We find a discontinuous transition in a short-range binary
lattice model with exactly two absorbing configurations. The dynamics combines the
self-complementary Day \& Night rule, voter copying with weight $p$, and a local minority
flip. At $p=1/2$, the active phase retains a finite interface density at the reversal of
a planar front. Bulk decay by nucleation and growth, front-orientation comparisons,
and droplets of each phase embedded in the other provide complementary evidence for
phase coexistence. A wall between the two absorbing phases develops an active layer
that thickens toward the transition. These observations distinguish a discontinuous
loss of activity from the magnetization jump that also occurs at a generalized voter
transition. Toward the voter point, the invasion boundary scales approximately as
$q^*\propto\alpha(1-p)$. A finite-time closure based on coalescing random walks accounts
for its scale and identifies contributions capable of bistability, while leaving the
asymptotic endpoint unresolved. The model provides a concrete setting for studying
nucleation and interfacial dynamics at a two-dimensional transition with two symmetric
absorbing configurations.
\end{abstract}
